\documentclass[11pt,letterpaper]{article}
\usepackage[T1]{fontenc}
\usepackage[utf8]{inputenc}
\usepackage{lmodern,microtype}
\usepackage[margin=0.78in,headheight=15pt]{geometry}
\usepackage{amsmath,amssymb,booktabs,longtable,tabularx,array,enumitem,xcolor,hyperref,fancyhdr,lastpage,needspace}
\definecolor{EngineBlue}{HTML}{183B56}
\definecolor{EngineTeal}{HTML}{117A83}
\definecolor{EngineGray}{HTML}{596671}
\definecolor{EngineLight}{HTML}{EEF4F7}
\hypersetup{colorlinks=true,urlcolor=EngineTeal,linkcolor=EngineBlue}
\setlength{\parindent}{0pt}
\setlength{\parskip}{5pt}
\emergencystretch=3em
\setlist{nosep,leftmargin=*}
\pagestyle{fancy}\fancyhf{}
\fancyhead[L]{\small\sffamily Game Engine Architecture | Engineering Suite}
\fancyhead[R]{\small\sffamily \DocType}
\fancyfoot[L]{\scriptsize Source-grounded engineering curriculum | v1.0}
\fancyfoot[R]{\scriptsize \thepage/\pageref{LastPage}}
\newcommand{\DocType}{Reference}
\newcommand{\SuiteHeading}[3]{%
\renewcommand{\DocType}{#2}%
{\small\bfseries\sffamily\color{EngineTeal}#2\hfill #3}\par\vspace{6pt}
{\LARGE\bfseries\color{EngineBlue}#1}\par\vspace{4pt}
{\color{EngineTeal}\rule{\linewidth}{1pt}}\par\vspace{7pt}}
\newcommand{\LabEvidence}[1]{\par\textbf{Evidence required:} #1\par}
\newcommand{\WriteLines}[1]{\par\begingroup\setlength{\parskip}{9pt}\count255=0\loop\ifnum\count255<#1\noindent\rule{\linewidth}{0.35pt}\par\advance\count255 by1\repeat\endgroup\smallskip}
\newcommand{\WritingBox}[1]{\noindent\fbox{\parbox[t][#1][t]{0.97\linewidth}{\mbox{}}}\par\smallskip}
\newcommand{\SourceNote}{\footnotesize\textit{Primary reading: Jason Gregory, \emph{Game Engine Architecture}, 3rd ed. (CRC Press, 2018). The narrative, exercises, laboratory designs and scoring instruments in this suite are independently authored teaching materials, not a reproduction of the textbook.}\normalsize\par}
\newcommand{\Criterion}[2]{\textbf{#1.} #2\par}
\newcommand{\FormField}[2]{\Needspace{3\baselineskip}\textbf{#1}\par\textcolor{EngineGray}{\footnotesize #2}\par\WriteLines{2}}
\newcommand{\BlankField}[1]{\Needspace{3\baselineskip}\textbf{#1}\par\WriteLines{2}}
\newcolumntype{Y}{>{\raggedright\arraybackslash}X}

\begin{document}
\SuiteHeading{Animation Runtime and State Machines}{Technical reference}{Runtime systems}
\SourceNote
\section{Scope and normative conventions}
This reference summarizes design options and independently proposed engineering controls. It is not a statement that the textbook mandates any particular API or modern implementation choice.

\subsection{Pose pipeline}
Clip time sampling yields local joint transforms; hierarchy propagation yields model-space transforms; skinning reads matrices or dual quaternions.

\begin{itemize}
\item Define the observable interface and resource owner.
\item Specify a success invariant and a deliberate failure case.
\item Trace at least one update frame from input to output.
\end{itemize}
\subsection{Blend correctness}
Blend translation/scale and rotations with appropriate methods; normalize quaternion interpolation and check hemisphere alignment to avoid long-arc turns.

\begin{itemize}
\item Define the observable interface and resource owner.
\item Specify a success invariant and a deliberate failure case.
\item Trace at least one update frame from input to output.
\end{itemize}
\subsection{State machine}
Transitions need conditions, interruption policy, duration, event delivery and root-motion arbitration; document whether the engine or gameplay owns transition choice.

\begin{itemize}
\item Define the observable interface and resource owner.
\item Specify a success invariant and a deliberate failure case.
\item Trace at least one update frame from input to output.
\end{itemize}
\subsection{Performance}
Track joint counts, clip memory, skinning submission and pose update cost; validate compressed clips against pose-error thresholds.

\begin{itemize}
\item Define the observable interface and resource owner.
\item Specify a success invariant and a deliberate failure case.
\item Trace at least one update frame from input to output.
\end{itemize}
\section{Design review questions}
\begin{itemize}
\item Where are time, memory and object lifetimes transferred across subsystem boundaries?
\item Which statements are assumptions, and what test would invalidate them?
\item What fails if content is unloaded while a worker thread still holds a reference?
\item Which metrics can establish that this subsystem satisfies a frame-time budget?
\end{itemize}
\section{Evidence package}
\begin{itemize}
\item Versioned contract and small dependency diagram.
\item Tests covering invalid inputs and state transitions.
\item Reproducible benchmark data, limitations and release decision.
\end{itemize}

\end{document}
